\chapter{Probability Current}

\section*{Introduction}

In quantum mechanics, particles do not have definite trajectories, so the notion of a ``current'' is subtle. The probability current $\mathbf{J}$ maps the abstract wavefunction dynamics onto a tangible flow: it tells us, for every point in space, how much probability is moving through that point and in what direction. It is the observable signature of quantum motion.
\section*{Fundamental Equation Used}

\begin{equation}
\mathbf{J} = \frac{\hbar}{2mi}\left(\psi^*\nabla\psi - \psi\nabla\psi^*\right)
\end{equation}
\section*{Assumptions}

\textbf{Assumptions:} Non-relativistic quantum mechanics; single-particle description; wavefunction $\psi$ is single-valued and normalizable.


\textbf{Limitations:} Does not apply to many-body systems without generalization; relativistic effects require the Dirac current.


\section*{Mathematical Derivation}

\textbf{Step 1: Start from the probability density and its time evolution.}

The probability density is $\rho(\mathbf{r},t) = |\psi|^2 = \psi^*\psi$. Its time derivative is
\begin{equation}
\frac{\partial\rho}{\partial t} = \frac{\partial\psi^*}{\partial t}\psi + \psi^*\frac{\partial\psi}{\partial t}.
\end{equation}

\textbf{Step 2: Substitute the Schr\"{o}dinger equation for $\partial\psi/\partial t$.}

The time-dependent Schr\"{o}dinger equation and its complex conjugate are
\begin{align}
i\hbar\frac{\partial\psi}{\partial t} &= -\frac{\hbar^2}{2m}\nabla^2\psi + V\psi, \\
-i\hbar\frac{\partial\psi^*}{\partial t} &= -\frac{\hbar^2}{2m}\nabla^2\psi^* + V\psi^*.
\end{align}
Solving for the time derivatives:
\begin{align}
\frac{\partial\psi}{\partial t} &= \frac{i\hbar}{2m}\nabla^2\psi - \frac{i}{\hbar}V\psi, \\
\frac{\partial\psi^*}{\partial t} &= -\frac{i\hbar}{2m}\nabla^2\psi^* + \frac{i}{\hbar}V\psi^*.
\end{align}

\textbf{Step 3: Insert into the time derivative of $\rho$.}

\begin{equation}
\frac{\partial\rho}{\partial t} = \left(-\frac{i\hbar}{2m}\nabla^2\psi^* + \frac{i}{\hbar}V\psi^*\right)\psi + \psi^*\left(\frac{i\hbar}{2m}\nabla^2\psi - \frac{i}{\hbar}V\psi\right).
\end{equation}
The potential terms cancel exactly:
\begin{equation}
\frac{\partial\rho}{\partial t} = \frac{i\hbar}{2m}\left(\psi^*\nabla^2\psi - \psi\nabla^2\psi^*\right).
\end{equation}

\textbf{Step 4: Rewrite as a divergence.}

Using the identity $\psi^*\nabla^2\psi - \psi\nabla^2\psi^* = \nabla\cdot(\psi^*\nabla\psi - \psi\nabla\psi^*)$, we obtain
\begin{equation}
\frac{\partial\rho}{\partial t} = \frac{i\hbar}{2m}\nabla\cdot\left(\psi^*\nabla\psi - \psi\nabla\psi^*\right).
\end{equation}

\textbf{Step 5: Define the probability current.}

Comparing with the continuity equation $\partial\rho/\partial t + \nabla\cdot\mathbf{J} = 0$, we identify
\begin{equation}
\boxed{\mathbf{J} = \frac{\hbar}{2mi}\left(\psi^*\nabla\psi - \psi\nabla\psi^*\right) = \frac{1}{2m}\left[\psi^*\hat{\mathbf{p}}\psi - \psi\hat{\mathbf{p}}\psi^*\right].}
\end{equation}
This is the probability current density. It is real (since $\mathbf{J}^* = \mathbf{J}$), it is gauge-covariant, and it reduces to $\rho\,\mathbf{v}$ in the classical limit.

\section*{Characteristics of the Equation}

\begin{itemize}
    \item $\mathbf{J}$ is a real vector field, ensuring physically meaningful probability flow.
    \item For a plane wave $\psi = e^{ikx}$, $\mathbf{J} = \frac{\hbar k}{m}|\psi|^2$, recovering the classical momentum density.
    \item The current vanishes for real-valued wavefunctions, reflecting stationary probability distributions.
    \item The probability current is gauge-invariant for the probability flow itself, though the canonical momentum is not.
\end{itemize}
\section*{Solution Methods}

\begin{itemize}
    \item \textbf{Direct computation:} Given $\psi$, compute $\nabla\psi$ and $\nabla\psi^*$ and combine.
    \item \textbf{Madelung formulation:} Write $\psi = \sqrt{\rho}\,e^{iS/\hbar}$ and identify $\mathbf{J} = \frac{\rho}{m}\nabla S$.
    \item \textbf{Operator method:} $\mathbf{J} = \frac{1}{2m}(\hat{\mathbf{p}}\delta(\mathbf{r}-\mathbf{r}') + \delta(\mathbf{r}-\mathbf{r}')\hat{\mathbf{p}})$ sandwiched between $\psi^*$ and $\psi$.
\end{itemize}
\section*{Verifications}

\begin{itemize}
    \item For a free-particle Gaussian wavepacket, the current correctly describes the spreading and center-of-mass motion.
    \item The integral of $\mathbf{J}$ over all space gives zero for bound states (no net probability flow out).
    \item The current reduces to $\rho \mathbf{v}$ in the classical limit (Ehrenfest theorem).
\end{itemize}
\section*{Applications}

\begin{itemize}
    \item Probability flow in scattering and tunneling problems.
    \item Persistent currents in mesoscopic rings and superconducting loops.
    \item Probability current in time-dependent potentials and driven quantum systems.
    \item Foundation for deriving the quantum continuity equation.
\end{itemize}
\section*{What Richard Feynman Would Have Asked}

\subsection*{Plain-English Translation}
\textit{``If a wavefunction is just a mathematical tool, what is the probability current actually telling us about the real physical world?''}

The probability current tells us the direction and rate at which the likelihood of finding a particle is flowing through space. It is the bridge between the abstract complex-valued wavefunction and the physical reality of particle motion. Even though we cannot track a quantum particle along a definite path, we can say with certainty how the probability distribution is reshaping itself over time, and the probability current is the quantity that encodes this reshaping.

\subsection*{Localized Mechanics}
\textit{``At a given point in space, what physical process does the probability current represent? Is there really something flowing there?''}

At any point, $\mathbf{J}$ represents the net rate at which probability crosses a tiny surface element at that point. It is not a flow of matter or energy per se---it is a flow of likelihood. If you placed an infinitesimal detector at that point, the probability current would tell you the net rate at which particle detections would drift past that detector. It is as real as the probability itself, which is to say it has direct physical consequences even though it describes something more abstract than classical fluid flow.

\subsection*{Testing Extreme Limits}
\textit{``What happens to the probability current when the potential is infinitely steep, like a perfect wall? Does current still flow?''}

At an impenetrable potential barrier, the probability current on the far side is exactly zero, and the reflected current equals the incident current on the near side. The current must be continuous at the boundary (even if the wavefunction has a kink), and this requirement determines the reflection coefficient. In the opposite limit of a completely flat potential, the current is uniform and constant, as expected for free propagation. These extreme cases confirm that the probability current correctly implements conservation of probability in all situations.

\subsection*{Dimensionality and Geometry}
\textit{``How does the probability current behave in two and three dimensions compared to one dimension? Is it fundamentally different?''}

In one dimension, $\mathbf{J}$ is a scalar with a sign indicating left or right flow. In higher dimensions, it becomes a vector field with richer structure---it can swirl, converge, or diverge. In two dimensions, probability currents can form vortices, which are directly related to orbital angular momentum in quantum systems. The three-dimensional case allows for the full complexity of probability flow seen in atomic orbitals and scattering problems. The dimensionality does not change the fundamental structure of the current formula, but it dramatically enriches the types of flow patterns that are possible.

\subsection*{Cross-Disciplinary Connection}
\textit{``Where in fluid mechanics, electromagnetism, or other fields do we encounter an analogous current concept, and how does it compare to the quantum probability current?''}

The probability current is the quantum analog of the mass flux $\rho\mathbf{v}$ in fluid dynamics and the electrical current density $\mathbf{J}$ in electromagnetism. In every case, the current is a flux---the product of a density and a velocity. The key difference is that in quantum mechanics, the velocity is not a property of the particle's state in the classical sense but is instead derived from the phase gradient of the wavefunction. This phase-velocity relationship is uniquely quantum: it connects the abstract complex phase of the wavefunction to the concrete physical flow of probability.

% ------------------------------------------------------------
\section*{FAQ}

\textbf{Q: Is probability current a measurable quantity?}
A: Not directly in the same way position or momentum are measured, but it has physical consequences. Interference patterns, persistent currents in mesoscopic systems, and persistent circulating currents in molecular orbitals all reflect the reality of probability flow.

\textbf{Q: Can probability current be negative?}
A: The components of $\mathbf{J}$ can be negative (indicating flow in the negative coordinate direction), but $\mathbf{J}$ is always real. There is no requirement for it to be positive-definite.

\textbf{Q: How does probability current relate to the electrical current?}
A: For a charged particle, the electrical current density is $\mathbf{j} = q\mathbf{J}$. The probability current is the universal kinematic object; the electrical current adds the charge factor.
\section*{Important Bibliography}

\begin{enumerate}
\item Griffiths, D.J., \textit{Introduction to Quantum Mechanics}, 3rd ed. (Cambridge University Press, 2018), Chapter 1.
\item Sakurai, J.J. and Napolitano, J., \textit{Modern Quantum Mechanics}, 2nd ed. (Cambridge University Press, 2017), Chapter 1.
\item Schiff, L.I., \textit{Quantum Mechanics}, 3rd ed. (McGraw-Hill, 1968), Chapter 1.
\item Feynman, R.P., \textit{The Feynman Lectures on Physics}, Vol.\ III (Addison-Wesley, 1965), Chapter 11.
\end{enumerate}


